\documentclass[a4paper]{report}

% Warn about commands, classes and packages that are outdated and superseded.
%     - Orthodox checks for pitfalls that are not technically incorrect.
\usepackage[l2tabu,orthodox,abort]{nag}

% All styling is contained in "mythesis.sty"
\usepackage{mythesis}

% Graphics directory
\graphicspath{{img/}}

% Title page options
\title{\thesisTitle*{}}
\author{\studentName*{}}
\date{\thesisDate*{}}

% Bibliography
\bibliography{thesis}

%%%%%%%%%%%%%%%%%%%%%%%%%%%%%%%%%%%%%%%%%%%%%%%%%%%%%%%%%%%%%%%%%%%%%%%%%%%%%%%%

\begin{document}

%%%%%%%%%%%%%%%%%%%%%%%%%%%%%%%%%%%%%%%%%%%%%%%%%
% Title Page
%%%%%%%%%%%%%%%%%%%%%%%%%%%%%%%%%%%%%%%%%%%%%%%%%
\begin{titlepage}
\centering

% Upper part of the page
\includegraphics[width=0.25\textwidth]{sydney_uni_coat_of_arms}\\[1cm]
\textsc{\LARGE \universityName*{}}\\[1.5cm]
\textsc{\Large Final year thesis}\\[0.5cm]

% Title
\rule{\linewidth}{0.5mm}\\[0.4cm] % horiizontal line
{\huge \bfseries \thesisTitle*{}}\\[0.4cm]
\rule{\linewidth}{0.5mm}\\[1.5cm] % horizontal line

% Author and supervisor
\begin{minipage}{0.4\textwidth}
    \begin{flushleft}
        \large
        \emph{Author:}\\
        \studentName*{}
    \end{flushleft}
\end{minipage}
\begin{minipage}{0.4\textwidth}
    \begin{flushright}
        \large
        \emph{Supervisor:}\\
        \supervisorName*{}
    \end{flushright}
\end{minipage}

% Strech vertical space so that it fills all empty space
\vfill

% Thesis description
{\large A thesis submitted in fulfilment of the requirements for the\\
degree of Bachelor of Engineering (Computer) in the\\
School of Electrical Engineering at\\
The University of Sydney}

% Strech vertical space so that it fills all empty space
\vfill

% Bottom of the page
{\large \thesisDate*{}}

\end{titlepage}

%%%%%%%%%%%%%%%%%%%%%%%%%%%%%%%%%%%%%%%%%%%%%%%%%
% Compliance Sheet
%%%%%%%%%%%%%%%%%%%%%%%%%%%%%%%%%%%%%%%%%%%%%%%%%
\includepdf[scale=0.8]{forms/compliance_sheet}

%%%%%%%%%%%%%%%%%%%%%%%%%%%%%%%%%%%%%%%%%%%%%%%%%
% Statement of Achievement
%%%%%%%%%%%%%%%%%%%%%%%%%%%%%%%%%%%%%%%%%%%%%%%%%
%\chapter*{Statement of Achievement}
% This is located in {chapter01/contributions}

\pagenumbering{roman}

%%%%%%%%%%%%%%%%%%%%%%%%%%%%%%%%%%%%%%%%%%%%%%%%%
% Acknowledgements
%%%%%%%%%%%%%%%%%%%%%%%%%%%%%%%%%%%%%%%%%%%%%%%%%
\chapter*{Acknowledgements}
This project could not have been completed without the guidance of my project 
supervisor, \supervisorName{}. He always made himself available to assist me 
with the understanding of many new and unknown concepts. His expertise in the 
field was a great asset and greatly contributed to my own understanding of the 
ideas explored in this thesis. His support and guidance was invaluable, and were
instrumental to the success of my project.

Much of the works completed during this project could not have been performed 
without the prior research of \getPerson{Chawla} and \getPerson{Khoa}. Both were
generous enough to lend their own support and expertise in order to gain an 
understanding of the implemented random projection algorithms.

Additionally, I must thank \getPerson{Frechtling} for his assistance and advice 
throughout the project. He assisted me with Xilnix tools such as ``AutoESL'', 
and also offered me valuable advice regarding my future professional career.

%%%%%%%%%%%%%%%%%%%%%%%%%%%%%%%%%%%%%%%%%%%%%%%%%
% Publications
%%%%%%%%%%%%%%%%%%%%%%%%%%%%%%%%%%%%%%%%%%%%%%%%%
\chapter*{Publications}
This thesis is based on the following publications:
\begin{enumerate}
\item \fullcite{Vries:2010}
\item \fullcite{Khoa:2012}
\item \fullcite{Bay:2003}
\end{enumerate}

%%%%%%%%%%%%%%%%%%%%%%%%%%%%%%%%%%%%%%%%%%%%%%%%%
% Abstract
%%%%%%%%%%%%%%%%%%%%%%%%%%%%%%%%%%%%%%%%%%%%%%%%%
\begin{abstract}
The prediction of the stock market has become an issue of great interest in the
areas of finance, mathematics and engineering; due mainly to the great potential
financial gain. Researchers have devised various algorithms and differing 
approaches to the problem of stock market analysis, with varying degrees of 
success.

A major outstanding issue for stock market analysis is the effective and 
efficient detection of local anomalies in the input data sets, which are 
inherently highly multidimensional. Many na\"{\i}ve algorithms are highly 
inefficient and others fail to adequately detect local anomalies altogether. It
had become a time-vs-correctness trade-off in which no acceptable compromise
could be reached.

However, researchers are starting to explore the relatively new concept of 
applying \doubleQuote{random projections} to the highly multidimensional data 
sets. Research has suggested that by applying these random projections, they are
able to significantly reduce the dimensionality (and consequently the
computational complexity) of the data sets, whilst sufficiently retaining the 
inherent properties of that data set --- at least so much so as anomaly
detection is concerned.

Anomaly detection is important because it allows otherwise-accurate machine 
learning algorithms such as neural networks to more accurately model and predict
the stock exchange data by ignoring anomalous data, which likely doesn't effect 
the state of the model to any significant degree.

\getPerson{Chawla} from the \universityName{} has in recent years conducted and 
supervised new and exciting research oriented around random projections. In 
particular, \getPerson{Khoa}, under the supervision of \getPerson{Chawla} 
evaluated the use of traditional distance metrics, such as Euclidean distance 
and Mahalanobis distance, in the application of local anomaly detection. 
\getPerson{Khoa} proposed the use of the \singleQuote{commute time} metric, 
derived from random walks on graphs, in anomaly detection.

The source code that was used throughout this project can be freely downloaded 
from \gitRepoHTTP{} (over HTTP) or \gitRepoSSH{} (over SSH).
\end{abstract}

%%%%%%%%%%%%%%%%%%%%%%%%%%%%%%%%%%%%%%%%%%%%%%%%%
% Table of Contents
%%%%%%%%%%%%%%%%%%%%%%%%%%%%%%%%%%%%%%%%%%%%%%%%%
\tableofcontents

%%%%%%%%%%%%%%%%%%%%%%%%%%%%%%%%%%%%%%%%%%%%%%%%%
% List of Figures
%%%%%%%%%%%%%%%%%%%%%%%%%%%%%%%%%%%%%%%%%%%%%%%%%
\listoffigures

%%%%%%%%%%%%%%%%%%%%%%%%%%%%%%%%%%%%%%%%%%%%%%%%%
% List of Algorithms
%%%%%%%%%%%%%%%%%%%%%%%%%%%%%%%%%%%%%%%%%%%%%%%%%
\listofalgorithms

%%%%%%%%%%%%%%%%%%%%%%%%%%%%%%%%%%%%%%%%%%%%%%%%%
% List of Tables
%%%%%%%%%%%%%%%%%%%%%%%%%%%%%%%%%%%%%%%%%%%%%%%%%
\listoftables

%%%%%%%%%%%%%%%%%%%%%%%%%%%%%%%%%%%%%%%%%%%%%%%%%
% Glossary
%%%%%%%%%%%%%%%%%%%%%%%%%%%%%%%%%%%%%%%%%%%%%%%%%
\glsaddall
\printglossary

\pagenumbering{arabic}

%%%%%%%%%%%%%%%%%%%%%%%%%%%%%%%%%%%%%%%%%%%%%%%%%
% CHAPTER 01: Introduction
%%%%%%%%%%%%%%%%%%%%%%%%%%%%%%%%%%%%%%%%%%%%%%%%%
\chapter{Introduction}
\label{introduction}

%%%%%%%%%%%%%%%%%%%%%%%%%%%%%%%%%%%%%%%%%%%%%%%%%
% Motivation
%%%%%%%%%%%%%%%%%%%%%%%%%%%%%%%%%%%%%%%%%%%%%%%%%
\section{Motivation}
\label{motivation}
Anomaly detection is an important technique which can be applied to a wide range
of applications. There are many techniques to detect and measure anomalies, each
usually most applicable to some specific problem domain \citeNeeded{}. A general
algorithm for anomaly detection has proved difficult to design \citeNeeded{},
due to various challenges associated with defining and measuring anomalies. Most
existing approaches are based on statistical or geometrical measures involving
distance. Whilst many algorithms work well for some subset of input data sets,
it has been difficult to discover an algorithm which performs both correctly and
efficiently on all possible data sets.

Previous research has found merit in applying randomization techniques to highly
multivariate data sets in order to reduce the dimensionality of these data sets
whilst maintaining their fundamental and statistical properties. Such reduction
of the dimensionality of data, assuming it can be performed efficiently, allows
previously-unscalable anomaly detection algorithms to be practically applied to
a wider range of data and applications.

Anomaly detection is an important and contemporary problem in the field of
computer science, and is of particular interesting to stock market analysis,
network intrusion detection and image comparison.

%%%%%%%%%%%%%%%%%%%%%%%%%%%%%%%%%%%%%%%%%%%%%%%%%
% Contributions of this thesis
%%%%%%%%%%%%%%%%%%%%%%%%%%%%%%%%%%%%%%%%%%%%%%%%%
\section{Contributions of this thesis}
\label{contributions}
One key difficulty in anomaly detection is the efficient scaling of a general
algorithm to apply to highly multivariate data \citeNeeded{}. In this thesis, I
explore the use of randomization techniques (such as random projections and
commute time) in anomaly detection algorithms. These techniques provide
encouraging results with regards to the run-time complexity of an algorithm.

Furthermore, in this thesis I attempt to make observations and an analysis of
the run-time performance of anomaly detection algorithms using randomization
techniques, so as to identify steps in the algorithms that bottleneck the
algorithm's performance. Through this identification it would be possible to
improve the \emph{actual} run-time performance of the algorithm by utilising
the advantages of reconfigurable computing.


%%%%%%%%%%%%%%%%%%%%%%%%%%%%%%%%%%%%%%%%%%%%%%%%%
% Organization
%%%%%%%%%%%%%%%%%%%%%%%%%%%%%%%%%%%%%%%%%%%%%%%%%
\section{Organization}
\label{organization}
The rest of this \thesis{} is organized as follows. In \autoref{background}, I
provide a background to various anomaly detection techniques, as well as
randomization techniques. In order to provide the reader with an understanding
of the background topics, a brief background is given to various topics in
linear algebra, vector calculus and graph theory. However, the provided
background material is not intended to be comprehensive, and it is expected that
the reader is familiar with these fields of mathematics. In
\autoref{reconfigurableComputing}, I provide an overview of reconfigurable
computing, including an explanation of \glspl{FPGA} and an overview of
\gls{HLS}.

In \autoref{software}, I profile the execution of an anomaly detection algorithm
and explore possible improvements to the algorithm, in particular by outsourcing
various stages of the algorithm to an FPGA device. I proceed to make iterative
refinements to the algorithm from its base implementation towards a hardware
design.

In \autoref{hardware}, I propose several hardware implementations and discuss
design considerations, as well as the design process itself, that must be
considered for these designs. I also detail the results of some initial testing
investigating the performance expectations from differing hardware designs.

I conclude in \autoref{conclusions} by reflecting upon the results obtained
through this research, and making suggestions for further research in this
topic.

 
%%%%%%%%%%%%%%%%%%%%%%%%%%%%%%%%%%%%%%%%%%%%%%%%%
% Schedule
%%%%%%%%%%%%%%%%%%%%%%%%%%%%%%%%%%%%%%%%%%%%%%%%%
\section{Schedule}
\label{schedule}
A Gantt chart showing the anticipated schedule for the project is shown in 
\autoref{fig:ganttChart}. This will be updated as the project progresses.

% Gantt chart
% See http://www.martin-kumm.de/tex_gantt_package.php
\begin{figure}[h]
	\centering
	\scalebox{0.5}{
	\begin{gantt}[
			xunitlength=0.5cm,
			fontsize=\small,
			titlefontsize=\small,
			drawledgerline=true
			]{18}{36}
		% Months
		\begin{ganttitle}
			\titleelement{Mar}{4}
			\titleelement{Apr}{4}
			\titleelement{May}{4}
			\titleelement{Jun}{4}
			\titleelement{Jul}{4}
			\titleelement{Aug}{4}
			\titleelement{Sep}{4}
			\titleelement{Oct}{4}
			\titleelement{Nov}{4}
		\end{ganttitle}
		% Split months into quarters
		\begin{ganttitle}
			\numtitle{1}{1}{4}{1}
			\numtitle{1}{1}{4}{1}
			\numtitle{1}{1}{4}{1}
			\numtitle{1}{1}{4}{1}
			\numtitle{1}{1}{4}{1}
			\numtitle{1}{1}{4}{1}
			\numtitle{1}{1}{4}{1}
			\numtitle{1}{1}{4}{1}
			\numtitle{1}{1}{4}{1}
		\end{ganttitle}

		\input{ganttdata}
	\end{gantt}
	}
	\caption{Schedule for thesis work}
	\label{fig:ganttChart}
\end{figure}



%%%%%%%%%%%%%%%%%%%%%%%%%%%%%%%%%%%%%%%%%%%%%%%%%
% CHAPTER 02: Background
%%%%%%%%%%%%%%%%%%%%%%%%%%%%%%%%%%%%%%%%%%%%%%%%%
\chapter{Introduction}
\label{introduction}

%%%%%%%%%%%%%%%%%%%%%%%%%%%%%%%%%%%%%%%%%%%%%%%%%
% Motivation
%%%%%%%%%%%%%%%%%%%%%%%%%%%%%%%%%%%%%%%%%%%%%%%%%
\section{Motivation}
\label{motivation}
Anomaly detection is an important technique which can be applied to a wide range
of applications. There are many techniques to detect and measure anomalies, each
usually most applicable to some specific problem domain \citeNeeded{}. A general
algorithm for anomaly detection has proved difficult to design \citeNeeded{},
due to various challenges associated with defining and measuring anomalies. Most
existing approaches are based on statistical or geometrical measures involving
distance. Whilst many algorithms work well for some subset of input data sets,
it has been difficult to discover an algorithm which performs both correctly and
efficiently on all possible data sets.

Previous research has found merit in applying randomization techniques to highly
multivariate data sets in order to reduce the dimensionality of these data sets
whilst maintaining their fundamental and statistical properties. Such reduction
of the dimensionality of data, assuming it can be performed efficiently, allows
previously-unscalable anomaly detection algorithms to be practically applied to
a wider range of data and applications.

Anomaly detection is an important and contemporary problem in the field of
computer science, and is of particular interesting to stock market analysis,
network intrusion detection and image comparison.

%%%%%%%%%%%%%%%%%%%%%%%%%%%%%%%%%%%%%%%%%%%%%%%%%
% Contributions of this thesis
%%%%%%%%%%%%%%%%%%%%%%%%%%%%%%%%%%%%%%%%%%%%%%%%%
\section{Contributions of this thesis}
\label{contributions}
One key difficulty in anomaly detection is the efficient scaling of a general
algorithm to apply to highly multivariate data \citeNeeded{}. In this thesis, I
explore the use of randomization techniques (such as random projections and
commute time) in anomaly detection algorithms. These techniques provide
encouraging results with regards to the run-time complexity of an algorithm.

Furthermore, in this thesis I attempt to make observations and an analysis of
the run-time performance of anomaly detection algorithms using randomization
techniques, so as to identify steps in the algorithms that bottleneck the
algorithm's performance. Through this identification it would be possible to
improve the \emph{actual} run-time performance of the algorithm by utilising
the advantages of reconfigurable computing.


%%%%%%%%%%%%%%%%%%%%%%%%%%%%%%%%%%%%%%%%%%%%%%%%%
% Organization
%%%%%%%%%%%%%%%%%%%%%%%%%%%%%%%%%%%%%%%%%%%%%%%%%
\section{Organization}
\label{organization}
The rest of this \thesis{} is organized as follows. In \autoref{background}, I
provide a background to various anomaly detection techniques, as well as
randomization techniques. In order to provide the reader with an understanding
of the background topics, a brief background is given to various topics in
linear algebra, vector calculus and graph theory. However, the provided
background material is not intended to be comprehensive, and it is expected that
the reader is familiar with these fields of mathematics. In
\autoref{reconfigurableComputing}, I provide an overview of reconfigurable
computing, including an explanation of \glspl{FPGA} and an overview of
\gls{HLS}.

In \autoref{software}, I profile the execution of an anomaly detection algorithm
and explore possible improvements to the algorithm, in particular by outsourcing
various stages of the algorithm to an FPGA device. I proceed to make iterative
refinements to the algorithm from its base implementation towards a hardware
design.

In \autoref{hardware}, I propose several hardware implementations and discuss
design considerations, as well as the design process itself, that must be
considered for these designs. I also detail the results of some initial testing
investigating the performance expectations from differing hardware designs.

I conclude in \autoref{conclusions} by reflecting upon the results obtained
through this research, and making suggestions for further research in this
topic.

 
%%%%%%%%%%%%%%%%%%%%%%%%%%%%%%%%%%%%%%%%%%%%%%%%%
% Schedule
%%%%%%%%%%%%%%%%%%%%%%%%%%%%%%%%%%%%%%%%%%%%%%%%%
\section{Schedule}
\label{schedule}
A Gantt chart showing the anticipated schedule for the project is shown in 
\autoref{fig:ganttChart}. This will be updated as the project progresses.

% Gantt chart
% See http://www.martin-kumm.de/tex_gantt_package.php
\begin{figure}[h]
	\centering
	\scalebox{0.5}{
	\begin{gantt}[
			xunitlength=0.5cm,
			fontsize=\small,
			titlefontsize=\small,
			drawledgerline=true
			]{18}{36}
		% Months
		\begin{ganttitle}
			\titleelement{Mar}{4}
			\titleelement{Apr}{4}
			\titleelement{May}{4}
			\titleelement{Jun}{4}
			\titleelement{Jul}{4}
			\titleelement{Aug}{4}
			\titleelement{Sep}{4}
			\titleelement{Oct}{4}
			\titleelement{Nov}{4}
		\end{ganttitle}
		% Split months into quarters
		\begin{ganttitle}
			\numtitle{1}{1}{4}{1}
			\numtitle{1}{1}{4}{1}
			\numtitle{1}{1}{4}{1}
			\numtitle{1}{1}{4}{1}
			\numtitle{1}{1}{4}{1}
			\numtitle{1}{1}{4}{1}
			\numtitle{1}{1}{4}{1}
			\numtitle{1}{1}{4}{1}
			\numtitle{1}{1}{4}{1}
		\end{ganttitle}

		\input{ganttdata}
	\end{gantt}
	}
	\caption{Schedule for thesis work}
	\label{fig:ganttChart}
\end{figure}



%%%%%%%%%%%%%%%%%%%%%%%%%%%%%%%%%%%%%%%%%%%%%%%%%
% CHAPTER 03: Reconfigurable Computing
%%%%%%%%%%%%%%%%%%%%%%%%%%%%%%%%%%%%%%%%%%%%%%%%%
\chapter{Introduction}
\label{introduction}

%%%%%%%%%%%%%%%%%%%%%%%%%%%%%%%%%%%%%%%%%%%%%%%%%
% Motivation
%%%%%%%%%%%%%%%%%%%%%%%%%%%%%%%%%%%%%%%%%%%%%%%%%
\section{Motivation}
\label{motivation}
Anomaly detection is an important technique which can be applied to a wide range
of applications. There are many techniques to detect and measure anomalies, each
usually most applicable to some specific problem domain \citeNeeded{}. A general
algorithm for anomaly detection has proved difficult to design \citeNeeded{},
due to various challenges associated with defining and measuring anomalies. Most
existing approaches are based on statistical or geometrical measures involving
distance. Whilst many algorithms work well for some subset of input data sets,
it has been difficult to discover an algorithm which performs both correctly and
efficiently on all possible data sets.

Previous research has found merit in applying randomization techniques to highly
multivariate data sets in order to reduce the dimensionality of these data sets
whilst maintaining their fundamental and statistical properties. Such reduction
of the dimensionality of data, assuming it can be performed efficiently, allows
previously-unscalable anomaly detection algorithms to be practically applied to
a wider range of data and applications.

Anomaly detection is an important and contemporary problem in the field of
computer science, and is of particular interesting to stock market analysis,
network intrusion detection and image comparison.

%%%%%%%%%%%%%%%%%%%%%%%%%%%%%%%%%%%%%%%%%%%%%%%%%
% Contributions of this thesis
%%%%%%%%%%%%%%%%%%%%%%%%%%%%%%%%%%%%%%%%%%%%%%%%%
\section{Contributions of this thesis}
\label{contributions}
One key difficulty in anomaly detection is the efficient scaling of a general
algorithm to apply to highly multivariate data \citeNeeded{}. In this thesis, I
explore the use of randomization techniques (such as random projections and
commute time) in anomaly detection algorithms. These techniques provide
encouraging results with regards to the run-time complexity of an algorithm.

Furthermore, in this thesis I attempt to make observations and an analysis of
the run-time performance of anomaly detection algorithms using randomization
techniques, so as to identify steps in the algorithms that bottleneck the
algorithm's performance. Through this identification it would be possible to
improve the \emph{actual} run-time performance of the algorithm by utilising
the advantages of reconfigurable computing.


%%%%%%%%%%%%%%%%%%%%%%%%%%%%%%%%%%%%%%%%%%%%%%%%%
% Organization
%%%%%%%%%%%%%%%%%%%%%%%%%%%%%%%%%%%%%%%%%%%%%%%%%
\section{Organization}
\label{organization}
The rest of this \thesis{} is organized as follows. In \autoref{background}, I
provide a background to various anomaly detection techniques, as well as
randomization techniques. In order to provide the reader with an understanding
of the background topics, a brief background is given to various topics in
linear algebra, vector calculus and graph theory. However, the provided
background material is not intended to be comprehensive, and it is expected that
the reader is familiar with these fields of mathematics. In
\autoref{reconfigurableComputing}, I provide an overview of reconfigurable
computing, including an explanation of \glspl{FPGA} and an overview of
\gls{HLS}.

In \autoref{software}, I profile the execution of an anomaly detection algorithm
and explore possible improvements to the algorithm, in particular by outsourcing
various stages of the algorithm to an FPGA device. I proceed to make iterative
refinements to the algorithm from its base implementation towards a hardware
design.

In \autoref{hardware}, I propose several hardware implementations and discuss
design considerations, as well as the design process itself, that must be
considered for these designs. I also detail the results of some initial testing
investigating the performance expectations from differing hardware designs.

I conclude in \autoref{conclusions} by reflecting upon the results obtained
through this research, and making suggestions for further research in this
topic.

 
%%%%%%%%%%%%%%%%%%%%%%%%%%%%%%%%%%%%%%%%%%%%%%%%%
% Schedule
%%%%%%%%%%%%%%%%%%%%%%%%%%%%%%%%%%%%%%%%%%%%%%%%%
\section{Schedule}
\label{schedule}
A Gantt chart showing the anticipated schedule for the project is shown in 
\autoref{fig:ganttChart}. This will be updated as the project progresses.

% Gantt chart
% See http://www.martin-kumm.de/tex_gantt_package.php
\begin{figure}[h]
	\centering
	\scalebox{0.5}{
	\begin{gantt}[
			xunitlength=0.5cm,
			fontsize=\small,
			titlefontsize=\small,
			drawledgerline=true
			]{18}{36}
		% Months
		\begin{ganttitle}
			\titleelement{Mar}{4}
			\titleelement{Apr}{4}
			\titleelement{May}{4}
			\titleelement{Jun}{4}
			\titleelement{Jul}{4}
			\titleelement{Aug}{4}
			\titleelement{Sep}{4}
			\titleelement{Oct}{4}
			\titleelement{Nov}{4}
		\end{ganttitle}
		% Split months into quarters
		\begin{ganttitle}
			\numtitle{1}{1}{4}{1}
			\numtitle{1}{1}{4}{1}
			\numtitle{1}{1}{4}{1}
			\numtitle{1}{1}{4}{1}
			\numtitle{1}{1}{4}{1}
			\numtitle{1}{1}{4}{1}
			\numtitle{1}{1}{4}{1}
			\numtitle{1}{1}{4}{1}
			\numtitle{1}{1}{4}{1}
		\end{ganttitle}

		\input{ganttdata}
	\end{gantt}
	}
	\caption{Schedule for thesis work}
	\label{fig:ganttChart}
\end{figure}



%%%%%%%%%%%%%%%%%%%%%%%%%%%%%%%%%%%%%%%%%%%%%%%%%
% CHAPTER 04: Design
%%%%%%%%%%%%%%%%%%%%%%%%%%%%%%%%%%%%%%%%%%%%%%%%%
\chapter{Introduction}
\label{introduction}

%%%%%%%%%%%%%%%%%%%%%%%%%%%%%%%%%%%%%%%%%%%%%%%%%
% Motivation
%%%%%%%%%%%%%%%%%%%%%%%%%%%%%%%%%%%%%%%%%%%%%%%%%
\section{Motivation}
\label{motivation}
Anomaly detection is an important technique which can be applied to a wide range
of applications. There are many techniques to detect and measure anomalies, each
usually most applicable to some specific problem domain \citeNeeded{}. A general
algorithm for anomaly detection has proved difficult to design \citeNeeded{},
due to various challenges associated with defining and measuring anomalies. Most
existing approaches are based on statistical or geometrical measures involving
distance. Whilst many algorithms work well for some subset of input data sets,
it has been difficult to discover an algorithm which performs both correctly and
efficiently on all possible data sets.

Previous research has found merit in applying randomization techniques to highly
multivariate data sets in order to reduce the dimensionality of these data sets
whilst maintaining their fundamental and statistical properties. Such reduction
of the dimensionality of data, assuming it can be performed efficiently, allows
previously-unscalable anomaly detection algorithms to be practically applied to
a wider range of data and applications.

Anomaly detection is an important and contemporary problem in the field of
computer science, and is of particular interesting to stock market analysis,
network intrusion detection and image comparison.

%%%%%%%%%%%%%%%%%%%%%%%%%%%%%%%%%%%%%%%%%%%%%%%%%
% Contributions of this thesis
%%%%%%%%%%%%%%%%%%%%%%%%%%%%%%%%%%%%%%%%%%%%%%%%%
\section{Contributions of this thesis}
\label{contributions}
One key difficulty in anomaly detection is the efficient scaling of a general
algorithm to apply to highly multivariate data \citeNeeded{}. In this thesis, I
explore the use of randomization techniques (such as random projections and
commute time) in anomaly detection algorithms. These techniques provide
encouraging results with regards to the run-time complexity of an algorithm.

Furthermore, in this thesis I attempt to make observations and an analysis of
the run-time performance of anomaly detection algorithms using randomization
techniques, so as to identify steps in the algorithms that bottleneck the
algorithm's performance. Through this identification it would be possible to
improve the \emph{actual} run-time performance of the algorithm by utilising
the advantages of reconfigurable computing.


%%%%%%%%%%%%%%%%%%%%%%%%%%%%%%%%%%%%%%%%%%%%%%%%%
% Organization
%%%%%%%%%%%%%%%%%%%%%%%%%%%%%%%%%%%%%%%%%%%%%%%%%
\section{Organization}
\label{organization}
The rest of this \thesis{} is organized as follows. In \autoref{background}, I
provide a background to various anomaly detection techniques, as well as
randomization techniques. In order to provide the reader with an understanding
of the background topics, a brief background is given to various topics in
linear algebra, vector calculus and graph theory. However, the provided
background material is not intended to be comprehensive, and it is expected that
the reader is familiar with these fields of mathematics. In
\autoref{reconfigurableComputing}, I provide an overview of reconfigurable
computing, including an explanation of \glspl{FPGA} and an overview of
\gls{HLS}.

In \autoref{software}, I profile the execution of an anomaly detection algorithm
and explore possible improvements to the algorithm, in particular by outsourcing
various stages of the algorithm to an FPGA device. I proceed to make iterative
refinements to the algorithm from its base implementation towards a hardware
design.

In \autoref{hardware}, I propose several hardware implementations and discuss
design considerations, as well as the design process itself, that must be
considered for these designs. I also detail the results of some initial testing
investigating the performance expectations from differing hardware designs.

I conclude in \autoref{conclusions} by reflecting upon the results obtained
through this research, and making suggestions for further research in this
topic.

 
%%%%%%%%%%%%%%%%%%%%%%%%%%%%%%%%%%%%%%%%%%%%%%%%%
% Schedule
%%%%%%%%%%%%%%%%%%%%%%%%%%%%%%%%%%%%%%%%%%%%%%%%%
\section{Schedule}
\label{schedule}
A Gantt chart showing the anticipated schedule for the project is shown in 
\autoref{fig:ganttChart}. This will be updated as the project progresses.

% Gantt chart
% See http://www.martin-kumm.de/tex_gantt_package.php
\begin{figure}[h]
	\centering
	\scalebox{0.5}{
	\begin{gantt}[
			xunitlength=0.5cm,
			fontsize=\small,
			titlefontsize=\small,
			drawledgerline=true
			]{18}{36}
		% Months
		\begin{ganttitle}
			\titleelement{Mar}{4}
			\titleelement{Apr}{4}
			\titleelement{May}{4}
			\titleelement{Jun}{4}
			\titleelement{Jul}{4}
			\titleelement{Aug}{4}
			\titleelement{Sep}{4}
			\titleelement{Oct}{4}
			\titleelement{Nov}{4}
		\end{ganttitle}
		% Split months into quarters
		\begin{ganttitle}
			\numtitle{1}{1}{4}{1}
			\numtitle{1}{1}{4}{1}
			\numtitle{1}{1}{4}{1}
			\numtitle{1}{1}{4}{1}
			\numtitle{1}{1}{4}{1}
			\numtitle{1}{1}{4}{1}
			\numtitle{1}{1}{4}{1}
			\numtitle{1}{1}{4}{1}
			\numtitle{1}{1}{4}{1}
		\end{ganttitle}

		\input{ganttdata}
	\end{gantt}
	}
	\caption{Schedule for thesis work}
	\label{fig:ganttChart}
\end{figure}



%%%%%%%%%%%%%%%%%%%%%%%%%%%%%%%%%%%%%%%%%%%%%%%%%
% CHAPTER 05: Implementation
%%%%%%%%%%%%%%%%%%%%%%%%%%%%%%%%%%%%%%%%%%%%%%%%%
\chapter{Introduction}
\label{introduction}

%%%%%%%%%%%%%%%%%%%%%%%%%%%%%%%%%%%%%%%%%%%%%%%%%
% Motivation
%%%%%%%%%%%%%%%%%%%%%%%%%%%%%%%%%%%%%%%%%%%%%%%%%
\section{Motivation}
\label{motivation}
Anomaly detection is an important technique which can be applied to a wide range
of applications. There are many techniques to detect and measure anomalies, each
usually most applicable to some specific problem domain \citeNeeded{}. A general
algorithm for anomaly detection has proved difficult to design \citeNeeded{},
due to various challenges associated with defining and measuring anomalies. Most
existing approaches are based on statistical or geometrical measures involving
distance. Whilst many algorithms work well for some subset of input data sets,
it has been difficult to discover an algorithm which performs both correctly and
efficiently on all possible data sets.

Previous research has found merit in applying randomization techniques to highly
multivariate data sets in order to reduce the dimensionality of these data sets
whilst maintaining their fundamental and statistical properties. Such reduction
of the dimensionality of data, assuming it can be performed efficiently, allows
previously-unscalable anomaly detection algorithms to be practically applied to
a wider range of data and applications.

Anomaly detection is an important and contemporary problem in the field of
computer science, and is of particular interesting to stock market analysis,
network intrusion detection and image comparison.

%%%%%%%%%%%%%%%%%%%%%%%%%%%%%%%%%%%%%%%%%%%%%%%%%
% Contributions of this thesis
%%%%%%%%%%%%%%%%%%%%%%%%%%%%%%%%%%%%%%%%%%%%%%%%%
\section{Contributions of this thesis}
\label{contributions}
One key difficulty in anomaly detection is the efficient scaling of a general
algorithm to apply to highly multivariate data \citeNeeded{}. In this thesis, I
explore the use of randomization techniques (such as random projections and
commute time) in anomaly detection algorithms. These techniques provide
encouraging results with regards to the run-time complexity of an algorithm.

Furthermore, in this thesis I attempt to make observations and an analysis of
the run-time performance of anomaly detection algorithms using randomization
techniques, so as to identify steps in the algorithms that bottleneck the
algorithm's performance. Through this identification it would be possible to
improve the \emph{actual} run-time performance of the algorithm by utilising
the advantages of reconfigurable computing.


%%%%%%%%%%%%%%%%%%%%%%%%%%%%%%%%%%%%%%%%%%%%%%%%%
% Organization
%%%%%%%%%%%%%%%%%%%%%%%%%%%%%%%%%%%%%%%%%%%%%%%%%
\section{Organization}
\label{organization}
The rest of this \thesis{} is organized as follows. In \autoref{background}, I
provide a background to various anomaly detection techniques, as well as
randomization techniques. In order to provide the reader with an understanding
of the background topics, a brief background is given to various topics in
linear algebra, vector calculus and graph theory. However, the provided
background material is not intended to be comprehensive, and it is expected that
the reader is familiar with these fields of mathematics. In
\autoref{reconfigurableComputing}, I provide an overview of reconfigurable
computing, including an explanation of \glspl{FPGA} and an overview of
\gls{HLS}.

In \autoref{software}, I profile the execution of an anomaly detection algorithm
and explore possible improvements to the algorithm, in particular by outsourcing
various stages of the algorithm to an FPGA device. I proceed to make iterative
refinements to the algorithm from its base implementation towards a hardware
design.

In \autoref{hardware}, I propose several hardware implementations and discuss
design considerations, as well as the design process itself, that must be
considered for these designs. I also detail the results of some initial testing
investigating the performance expectations from differing hardware designs.

I conclude in \autoref{conclusions} by reflecting upon the results obtained
through this research, and making suggestions for further research in this
topic.

 
%%%%%%%%%%%%%%%%%%%%%%%%%%%%%%%%%%%%%%%%%%%%%%%%%
% Schedule
%%%%%%%%%%%%%%%%%%%%%%%%%%%%%%%%%%%%%%%%%%%%%%%%%
\section{Schedule}
\label{schedule}
A Gantt chart showing the anticipated schedule for the project is shown in 
\autoref{fig:ganttChart}. This will be updated as the project progresses.

% Gantt chart
% See http://www.martin-kumm.de/tex_gantt_package.php
\begin{figure}[h]
	\centering
	\scalebox{0.5}{
	\begin{gantt}[
			xunitlength=0.5cm,
			fontsize=\small,
			titlefontsize=\small,
			drawledgerline=true
			]{18}{36}
		% Months
		\begin{ganttitle}
			\titleelement{Mar}{4}
			\titleelement{Apr}{4}
			\titleelement{May}{4}
			\titleelement{Jun}{4}
			\titleelement{Jul}{4}
			\titleelement{Aug}{4}
			\titleelement{Sep}{4}
			\titleelement{Oct}{4}
			\titleelement{Nov}{4}
		\end{ganttitle}
		% Split months into quarters
		\begin{ganttitle}
			\numtitle{1}{1}{4}{1}
			\numtitle{1}{1}{4}{1}
			\numtitle{1}{1}{4}{1}
			\numtitle{1}{1}{4}{1}
			\numtitle{1}{1}{4}{1}
			\numtitle{1}{1}{4}{1}
			\numtitle{1}{1}{4}{1}
			\numtitle{1}{1}{4}{1}
			\numtitle{1}{1}{4}{1}
		\end{ganttitle}

		\input{ganttdata}
	\end{gantt}
	}
	\caption{Schedule for thesis work}
	\label{fig:ganttChart}
\end{figure}



%%%%%%%%%%%%%%%%%%%%%%%%%%%%%%%%%%%%%%%%%%%%%%%%%
% CHAPTER 06: Results
%%%%%%%%%%%%%%%%%%%%%%%%%%%%%%%%%%%%%%%%%%%%%%%%%
\chapter{Introduction}
\label{introduction}

%%%%%%%%%%%%%%%%%%%%%%%%%%%%%%%%%%%%%%%%%%%%%%%%%
% Motivation
%%%%%%%%%%%%%%%%%%%%%%%%%%%%%%%%%%%%%%%%%%%%%%%%%
\section{Motivation}
\label{motivation}
Anomaly detection is an important technique which can be applied to a wide range
of applications. There are many techniques to detect and measure anomalies, each
usually most applicable to some specific problem domain \citeNeeded{}. A general
algorithm for anomaly detection has proved difficult to design \citeNeeded{},
due to various challenges associated with defining and measuring anomalies. Most
existing approaches are based on statistical or geometrical measures involving
distance. Whilst many algorithms work well for some subset of input data sets,
it has been difficult to discover an algorithm which performs both correctly and
efficiently on all possible data sets.

Previous research has found merit in applying randomization techniques to highly
multivariate data sets in order to reduce the dimensionality of these data sets
whilst maintaining their fundamental and statistical properties. Such reduction
of the dimensionality of data, assuming it can be performed efficiently, allows
previously-unscalable anomaly detection algorithms to be practically applied to
a wider range of data and applications.

Anomaly detection is an important and contemporary problem in the field of
computer science, and is of particular interesting to stock market analysis,
network intrusion detection and image comparison.

%%%%%%%%%%%%%%%%%%%%%%%%%%%%%%%%%%%%%%%%%%%%%%%%%
% Contributions of this thesis
%%%%%%%%%%%%%%%%%%%%%%%%%%%%%%%%%%%%%%%%%%%%%%%%%
\section{Contributions of this thesis}
\label{contributions}
One key difficulty in anomaly detection is the efficient scaling of a general
algorithm to apply to highly multivariate data \citeNeeded{}. In this thesis, I
explore the use of randomization techniques (such as random projections and
commute time) in anomaly detection algorithms. These techniques provide
encouraging results with regards to the run-time complexity of an algorithm.

Furthermore, in this thesis I attempt to make observations and an analysis of
the run-time performance of anomaly detection algorithms using randomization
techniques, so as to identify steps in the algorithms that bottleneck the
algorithm's performance. Through this identification it would be possible to
improve the \emph{actual} run-time performance of the algorithm by utilising
the advantages of reconfigurable computing.


%%%%%%%%%%%%%%%%%%%%%%%%%%%%%%%%%%%%%%%%%%%%%%%%%
% Organization
%%%%%%%%%%%%%%%%%%%%%%%%%%%%%%%%%%%%%%%%%%%%%%%%%
\section{Organization}
\label{organization}
The rest of this \thesis{} is organized as follows. In \autoref{background}, I
provide a background to various anomaly detection techniques, as well as
randomization techniques. In order to provide the reader with an understanding
of the background topics, a brief background is given to various topics in
linear algebra, vector calculus and graph theory. However, the provided
background material is not intended to be comprehensive, and it is expected that
the reader is familiar with these fields of mathematics. In
\autoref{reconfigurableComputing}, I provide an overview of reconfigurable
computing, including an explanation of \glspl{FPGA} and an overview of
\gls{HLS}.

In \autoref{software}, I profile the execution of an anomaly detection algorithm
and explore possible improvements to the algorithm, in particular by outsourcing
various stages of the algorithm to an FPGA device. I proceed to make iterative
refinements to the algorithm from its base implementation towards a hardware
design.

In \autoref{hardware}, I propose several hardware implementations and discuss
design considerations, as well as the design process itself, that must be
considered for these designs. I also detail the results of some initial testing
investigating the performance expectations from differing hardware designs.

I conclude in \autoref{conclusions} by reflecting upon the results obtained
through this research, and making suggestions for further research in this
topic.

 
%%%%%%%%%%%%%%%%%%%%%%%%%%%%%%%%%%%%%%%%%%%%%%%%%
% Schedule
%%%%%%%%%%%%%%%%%%%%%%%%%%%%%%%%%%%%%%%%%%%%%%%%%
\section{Schedule}
\label{schedule}
A Gantt chart showing the anticipated schedule for the project is shown in 
\autoref{fig:ganttChart}. This will be updated as the project progresses.

% Gantt chart
% See http://www.martin-kumm.de/tex_gantt_package.php
\begin{figure}[h]
	\centering
	\scalebox{0.5}{
	\begin{gantt}[
			xunitlength=0.5cm,
			fontsize=\small,
			titlefontsize=\small,
			drawledgerline=true
			]{18}{36}
		% Months
		\begin{ganttitle}
			\titleelement{Mar}{4}
			\titleelement{Apr}{4}
			\titleelement{May}{4}
			\titleelement{Jun}{4}
			\titleelement{Jul}{4}
			\titleelement{Aug}{4}
			\titleelement{Sep}{4}
			\titleelement{Oct}{4}
			\titleelement{Nov}{4}
		\end{ganttitle}
		% Split months into quarters
		\begin{ganttitle}
			\numtitle{1}{1}{4}{1}
			\numtitle{1}{1}{4}{1}
			\numtitle{1}{1}{4}{1}
			\numtitle{1}{1}{4}{1}
			\numtitle{1}{1}{4}{1}
			\numtitle{1}{1}{4}{1}
			\numtitle{1}{1}{4}{1}
			\numtitle{1}{1}{4}{1}
			\numtitle{1}{1}{4}{1}
		\end{ganttitle}

		\input{ganttdata}
	\end{gantt}
	}
	\caption{Schedule for thesis work}
	\label{fig:ganttChart}
\end{figure}



%%%%%%%%%%%%%%%%%%%%%%%%%%%%%%%%%%%%%%%%%%%%%%%%%
% CHAPTER 07: Conclusions
%%%%%%%%%%%%%%%%%%%%%%%%%%%%%%%%%%%%%%%%%%%%%%%%%
\chapter{Introduction}
\label{introduction}

%%%%%%%%%%%%%%%%%%%%%%%%%%%%%%%%%%%%%%%%%%%%%%%%%
% Motivation
%%%%%%%%%%%%%%%%%%%%%%%%%%%%%%%%%%%%%%%%%%%%%%%%%
\section{Motivation}
\label{motivation}
Anomaly detection is an important technique which can be applied to a wide range
of applications. There are many techniques to detect and measure anomalies, each
usually most applicable to some specific problem domain \citeNeeded{}. A general
algorithm for anomaly detection has proved difficult to design \citeNeeded{},
due to various challenges associated with defining and measuring anomalies. Most
existing approaches are based on statistical or geometrical measures involving
distance. Whilst many algorithms work well for some subset of input data sets,
it has been difficult to discover an algorithm which performs both correctly and
efficiently on all possible data sets.

Previous research has found merit in applying randomization techniques to highly
multivariate data sets in order to reduce the dimensionality of these data sets
whilst maintaining their fundamental and statistical properties. Such reduction
of the dimensionality of data, assuming it can be performed efficiently, allows
previously-unscalable anomaly detection algorithms to be practically applied to
a wider range of data and applications.

Anomaly detection is an important and contemporary problem in the field of
computer science, and is of particular interesting to stock market analysis,
network intrusion detection and image comparison.

%%%%%%%%%%%%%%%%%%%%%%%%%%%%%%%%%%%%%%%%%%%%%%%%%
% Contributions of this thesis
%%%%%%%%%%%%%%%%%%%%%%%%%%%%%%%%%%%%%%%%%%%%%%%%%
\section{Contributions of this thesis}
\label{contributions}
One key difficulty in anomaly detection is the efficient scaling of a general
algorithm to apply to highly multivariate data \citeNeeded{}. In this thesis, I
explore the use of randomization techniques (such as random projections and
commute time) in anomaly detection algorithms. These techniques provide
encouraging results with regards to the run-time complexity of an algorithm.

Furthermore, in this thesis I attempt to make observations and an analysis of
the run-time performance of anomaly detection algorithms using randomization
techniques, so as to identify steps in the algorithms that bottleneck the
algorithm's performance. Through this identification it would be possible to
improve the \emph{actual} run-time performance of the algorithm by utilising
the advantages of reconfigurable computing.


%%%%%%%%%%%%%%%%%%%%%%%%%%%%%%%%%%%%%%%%%%%%%%%%%
% Organization
%%%%%%%%%%%%%%%%%%%%%%%%%%%%%%%%%%%%%%%%%%%%%%%%%
\section{Organization}
\label{organization}
The rest of this \thesis{} is organized as follows. In \autoref{background}, I
provide a background to various anomaly detection techniques, as well as
randomization techniques. In order to provide the reader with an understanding
of the background topics, a brief background is given to various topics in
linear algebra, vector calculus and graph theory. However, the provided
background material is not intended to be comprehensive, and it is expected that
the reader is familiar with these fields of mathematics. In
\autoref{reconfigurableComputing}, I provide an overview of reconfigurable
computing, including an explanation of \glspl{FPGA} and an overview of
\gls{HLS}.

In \autoref{software}, I profile the execution of an anomaly detection algorithm
and explore possible improvements to the algorithm, in particular by outsourcing
various stages of the algorithm to an FPGA device. I proceed to make iterative
refinements to the algorithm from its base implementation towards a hardware
design.

In \autoref{hardware}, I propose several hardware implementations and discuss
design considerations, as well as the design process itself, that must be
considered for these designs. I also detail the results of some initial testing
investigating the performance expectations from differing hardware designs.

I conclude in \autoref{conclusions} by reflecting upon the results obtained
through this research, and making suggestions for further research in this
topic.

 
%%%%%%%%%%%%%%%%%%%%%%%%%%%%%%%%%%%%%%%%%%%%%%%%%
% Schedule
%%%%%%%%%%%%%%%%%%%%%%%%%%%%%%%%%%%%%%%%%%%%%%%%%
\section{Schedule}
\label{schedule}
A Gantt chart showing the anticipated schedule for the project is shown in 
\autoref{fig:ganttChart}. This will be updated as the project progresses.

% Gantt chart
% See http://www.martin-kumm.de/tex_gantt_package.php
\begin{figure}[h]
	\centering
	\scalebox{0.5}{
	\begin{gantt}[
			xunitlength=0.5cm,
			fontsize=\small,
			titlefontsize=\small,
			drawledgerline=true
			]{18}{36}
		% Months
		\begin{ganttitle}
			\titleelement{Mar}{4}
			\titleelement{Apr}{4}
			\titleelement{May}{4}
			\titleelement{Jun}{4}
			\titleelement{Jul}{4}
			\titleelement{Aug}{4}
			\titleelement{Sep}{4}
			\titleelement{Oct}{4}
			\titleelement{Nov}{4}
		\end{ganttitle}
		% Split months into quarters
		\begin{ganttitle}
			\numtitle{1}{1}{4}{1}
			\numtitle{1}{1}{4}{1}
			\numtitle{1}{1}{4}{1}
			\numtitle{1}{1}{4}{1}
			\numtitle{1}{1}{4}{1}
			\numtitle{1}{1}{4}{1}
			\numtitle{1}{1}{4}{1}
			\numtitle{1}{1}{4}{1}
			\numtitle{1}{1}{4}{1}
		\end{ganttitle}

		\input{ganttdata}
	\end{gantt}
	}
	\caption{Schedule for thesis work}
	\label{fig:ganttChart}
\end{figure}



%%%%%%%%%%%%%%%%%%%%%%%%%%%%%%%%%%%%%%%%%%%%%%%%%
% Appendix
%%%%%%%%%%%%%%%%%%%%%%%%%%%%%%%%%%%%%%%%%%%%%%%%%
\chapter{Introduction}
\label{introduction}

%%%%%%%%%%%%%%%%%%%%%%%%%%%%%%%%%%%%%%%%%%%%%%%%%
% Motivation
%%%%%%%%%%%%%%%%%%%%%%%%%%%%%%%%%%%%%%%%%%%%%%%%%
\section{Motivation}
\label{motivation}
Anomaly detection is an important technique which can be applied to a wide range
of applications. There are many techniques to detect and measure anomalies, each
usually most applicable to some specific problem domain \citeNeeded{}. A general
algorithm for anomaly detection has proved difficult to design \citeNeeded{},
due to various challenges associated with defining and measuring anomalies. Most
existing approaches are based on statistical or geometrical measures involving
distance. Whilst many algorithms work well for some subset of input data sets,
it has been difficult to discover an algorithm which performs both correctly and
efficiently on all possible data sets.

Previous research has found merit in applying randomization techniques to highly
multivariate data sets in order to reduce the dimensionality of these data sets
whilst maintaining their fundamental and statistical properties. Such reduction
of the dimensionality of data, assuming it can be performed efficiently, allows
previously-unscalable anomaly detection algorithms to be practically applied to
a wider range of data and applications.

Anomaly detection is an important and contemporary problem in the field of
computer science, and is of particular interesting to stock market analysis,
network intrusion detection and image comparison.

%%%%%%%%%%%%%%%%%%%%%%%%%%%%%%%%%%%%%%%%%%%%%%%%%
% Contributions of this thesis
%%%%%%%%%%%%%%%%%%%%%%%%%%%%%%%%%%%%%%%%%%%%%%%%%
\section{Contributions of this thesis}
\label{contributions}
One key difficulty in anomaly detection is the efficient scaling of a general
algorithm to apply to highly multivariate data \citeNeeded{}. In this thesis, I
explore the use of randomization techniques (such as random projections and
commute time) in anomaly detection algorithms. These techniques provide
encouraging results with regards to the run-time complexity of an algorithm.

Furthermore, in this thesis I attempt to make observations and an analysis of
the run-time performance of anomaly detection algorithms using randomization
techniques, so as to identify steps in the algorithms that bottleneck the
algorithm's performance. Through this identification it would be possible to
improve the \emph{actual} run-time performance of the algorithm by utilising
the advantages of reconfigurable computing.


%%%%%%%%%%%%%%%%%%%%%%%%%%%%%%%%%%%%%%%%%%%%%%%%%
% Organization
%%%%%%%%%%%%%%%%%%%%%%%%%%%%%%%%%%%%%%%%%%%%%%%%%
\section{Organization}
\label{organization}
The rest of this \thesis{} is organized as follows. In \autoref{background}, I
provide a background to various anomaly detection techniques, as well as
randomization techniques. In order to provide the reader with an understanding
of the background topics, a brief background is given to various topics in
linear algebra, vector calculus and graph theory. However, the provided
background material is not intended to be comprehensive, and it is expected that
the reader is familiar with these fields of mathematics. In
\autoref{reconfigurableComputing}, I provide an overview of reconfigurable
computing, including an explanation of \glspl{FPGA} and an overview of
\gls{HLS}.

In \autoref{software}, I profile the execution of an anomaly detection algorithm
and explore possible improvements to the algorithm, in particular by outsourcing
various stages of the algorithm to an FPGA device. I proceed to make iterative
refinements to the algorithm from its base implementation towards a hardware
design.

In \autoref{hardware}, I propose several hardware implementations and discuss
design considerations, as well as the design process itself, that must be
considered for these designs. I also detail the results of some initial testing
investigating the performance expectations from differing hardware designs.

I conclude in \autoref{conclusions} by reflecting upon the results obtained
through this research, and making suggestions for further research in this
topic.

 
%%%%%%%%%%%%%%%%%%%%%%%%%%%%%%%%%%%%%%%%%%%%%%%%%
% Schedule
%%%%%%%%%%%%%%%%%%%%%%%%%%%%%%%%%%%%%%%%%%%%%%%%%
\section{Schedule}
\label{schedule}
A Gantt chart showing the anticipated schedule for the project is shown in 
\autoref{fig:ganttChart}. This will be updated as the project progresses.

% Gantt chart
% See http://www.martin-kumm.de/tex_gantt_package.php
\begin{figure}[h]
	\centering
	\scalebox{0.5}{
	\begin{gantt}[
			xunitlength=0.5cm,
			fontsize=\small,
			titlefontsize=\small,
			drawledgerline=true
			]{18}{36}
		% Months
		\begin{ganttitle}
			\titleelement{Mar}{4}
			\titleelement{Apr}{4}
			\titleelement{May}{4}
			\titleelement{Jun}{4}
			\titleelement{Jul}{4}
			\titleelement{Aug}{4}
			\titleelement{Sep}{4}
			\titleelement{Oct}{4}
			\titleelement{Nov}{4}
		\end{ganttitle}
		% Split months into quarters
		\begin{ganttitle}
			\numtitle{1}{1}{4}{1}
			\numtitle{1}{1}{4}{1}
			\numtitle{1}{1}{4}{1}
			\numtitle{1}{1}{4}{1}
			\numtitle{1}{1}{4}{1}
			\numtitle{1}{1}{4}{1}
			\numtitle{1}{1}{4}{1}
			\numtitle{1}{1}{4}{1}
			\numtitle{1}{1}{4}{1}
		\end{ganttitle}

		\input{ganttdata}
	\end{gantt}
	}
	\caption{Schedule for thesis work}
	\label{fig:ganttChart}
\end{figure}



%%%%%%%%%%%%%%%%%%%%%%%%%%%%%%%%%%%%%%%%%%%%%%%%%
% Bibliography
%%%%%%%%%%%%%%%%%%%%%%%%%%%%%%%%%%%%%%%%%%%%%%%%%
\nocite{*}
\printbibliography

\end{document}
